% ============================================================================
%  2  Grassmann 流形的几何结构
% ============================================================================
\section{Grassmann 流形的几何结构}
\label{sec:geometry}

\subsection{定义与典范丛}

记 $\GK=G_k(\R^{k+n})$，即 $\R^{k+n}$ 中所有无向 $k$ 维线性子空间组成的
Grassmann 流形；其上的点 $V\in\GK$ 本身就是一个 $k$ 维子空间。典范丛
$\cG_k\to\GK$ 定义为
\begin{equation}
  (\cG_k)_V=V .
  \label{eq:tautological}
\end{equation}

选定标准内积后，$V$ 有正交补 $V^{\perp}$，由此得到秩为 $n$ 的正交补丛
$\cGc_k$，其纤维为 $(\cGc_k)_V=V^{\perp}$。于是每个纤维都有正交直和分解
$V\oplus V^{\perp}=\R^{k+n}$，且随 $V$ 光滑变化，从而得到稳定平凡化
\begin{equation}
  \cG_k\oplus\cGc_k\cong\eps^{k+n}.
  \label{eq:splitting}
\end{equation}
这一关系是后续全部特征类计算的起点。若记 $w_i=w_i(\cG_k)$，则由 Whitney
乘积公式可得
\begin{equation}
  w(\cG_k)\,w(\cGc_k)=1,
  \qquad\text{即}\qquad
  w(\cGc_k)=w(\cG_k)^{-1}.
  \label{eq:whitney-inverse}
\end{equation}
这里的逆元是在 $H^*(\GK;\Ftwo)$ 的增次理想中取形式逆，并最终在 Grassmann
流形的真实上同调环中加以解释。

\subsection{正交补丛与稳定平凡化}

正交补丛并不是切丛。前者的秩为 $n$，而后者的秩为 $kn$；二者通过 $\Hom$ 丛相
联系。特别地，$w_i(\cGc_k)$ 与切丛的共轭类 $\wbar_i(T\GK)$ 一般并不是同一个类，
因此不能共用同一条横线记号。

\subsection{切丛公式}

\begin{proposition}[Grassmann 流形的切丛公式]
\label{prop:tangent}
存在自然的实向量丛同构
\begin{equation}
  T\GK\cong\Hom(\cG_k,\cGc_k)\cong\cG_k^{*}\otimes\cGc_k .
  \label{eq:tangent}
\end{equation}
\end{proposition}

\begin{proof}[证明（个人复现）]
固定 $V\in\GK$，取 $V^{\perp}$ 作为补空间。对充分小的线性映射
$A\colon V\to V^{\perp}$，令
\[
  \Gamma_A=\{\,v+A(v):v\in V\,\}.
\]

\begin{center}
\begin{minipage}{0.94\textwidth}
\centering
\begin{tikzpicture}[scale=1.15,>=Stealth,line join=round,line cap=round]
  % V：水平面（图中以其截线表示）
  \fill[black!6] (-0.3,-0.13) rectangle (6.3,0.13);
  \draw[thick] (0,0) -- (6.1,0) node[right] {$V$};
  % V^\perp：竖直方向
  \draw[thick,->] (0,0) -- (0,3.3) node[above] {$V^{\perp}$};
  % 原点
  \fill (0,0) circle (1.3pt) node[below left] {$0$};
  % Gamma_A：A 的图像
  \draw[thick] (0,0) -- (5.7,2.55) node[right] {$\Gamma_A$};
  % 取一点 v
  \coordinate (vb) at (3.5,0);
  \coordinate (vt) at (3.5,1.565);
  \draw[dashed,thin] (0,1.565) -- (vt);
  \draw[->,very thick] (vb) -- (vt) node[midway,left=4pt] {$A(v)$};
  \fill (vb) circle (1.4pt) node[below] {$v$};
  \fill (vt) circle (1.4pt) node[below right=1pt] {$v+A(v)$};
\end{tikzpicture}
\captionof{figure}{图像坐标 $A\mapsto\Gamma_A$。$V$ 为水平面（图中取其截线），
$V^{\perp}$ 为竖直方向，$\Gamma_A$ 是 $A$ 的图像。$A(v)$ 是点 $v$ 到
$\Gamma_A$ 的竖直位移；$A\mapsto\Gamma_A$ 因此给出
$T_V\GK\cong\Hom(V,V^{\perp})$。}
\label{fig:graph-coords}
\end{minipage}
\end{center}
则 $A\mapsto\Gamma_A$ 是 $V$ 的邻域的坐标，故
$T_V\GK\cong\Hom(V,V^{\perp})$。这些纤维同构随 $V$ 光滑变化，得到第一式；第二式
是 $\Hom(E,F)\cong E^{*}\otimes F$ 的自然同构。

Hsiang--Szczarba \cite{HsiangSzczarba1964} 的原始证明给出了同一公式的稳定版本：
若 $\xi$ 为典范丛，则
\[
  \tau(G_{N-k,k})\oplus(\xi\otimes\xi^{*})\cong N\xi_{\mathbb{R}} .
\]
在实情形用度量识别 $\xi^{*}\cong\xi$，置 $N=n+k$ 即得
\begin{equation}
  T\GK\oplus(\cG_k\otimes\cG_k)\cong(n+k)\,\cG_k .
  \label{eq:stable-tangent}
\end{equation}
这也可由 $\cG_k\oplus\cGc_k\cong\eps^{n+k}$ 与
$T\GK\cong\cG_k^{*}\otimes\cGc_k$ 直接推出。
\end{proof}

\begin{remark}
关系 \eqref{eq:stable-tangent} 是第 \ref{sec:classical} 章一切特征类计算的出发点；
它与命题 \ref{prop:tangent} 的局部形式互为补充：前者稳定、便于求总类，后者逐点、
便于看清 $\cGc_k$ 的几何来源。
\end{remark}
